\section{值域}

\begin{definition}[值域]
    函数 $f(x)$ 的值域是指 $f(x)$ 可能取到的所有值的集合，即 $y$ 的取值范围。
\end{definition}

\subsection{二次比二次}

\begin{definition}
    若一个二次函数比二次函数，有特殊求法和通解求法.
\end{definition}

\subsubsection{简单的二次比二次}

\subsubsection{完整的二次比二次}

\begin{theorem}[完整的二次比二次]
    形如：$\frac{a_1x^2 + b_1x + c_1}{a_2x^2 + b_2x + c_2}$ 的二次函数的求法考虑两种：\textbf{求导法} 和 \textbf{判别式法}
\end{theorem}

\begin{example}[求导法]
    求函数 $y = \frac{(1+x^2)^2}{(3+4x^2)(3x^2+4)}$ 的值域.
\end{example}

\begin{solution}
    令 $x^2=t$, 则：
    $$y=\frac{(1+t)^2}{(3+4 t)(3 t+4)}=\frac{t^2+2 t+1}{12 t^2+25 t+12}$$
\end{solution}

\begin{example}[判别式法]
    求函数 $y = \frac{(1+x^2)^2}{(3+4x^2)(3x^2+4)}$ 的值域.
\end{example}

\begin{solution}

\end{solution}

\subsubsection{炫技式的均值不等式求法}

\begin{example}[均值不等式求法]
    求函数 $y = \frac{(1+x^2)^2}{(3+4x^2)(3x^2+4)}$ 的值域.
\end{example}

\begin{solution}
    $$
        \frac{\left(1+x^2\right)^2}{\left(3+4 x^2\right)\left(3 x^2+4\right)} \leqslant \frac{\left(1+x^2\right)^2}{\left(\frac{7+7 x^2}{2}\right)^2}=\frac{4}{49} .
    $$
\end{solution}

